%=============================================================================!
% 1.10  WHAT WE NOW HAVE                                                      !
%=============================================================================!
\unnumbered{What we now have}

Position, velocity and acceleration are successive derivatives with
respect to time, and a motion follows from its acceleration once the
initial position and velocity are known. The sine and cosine are the
coordinates of a point on the unit circle, and with angles measured in
radians their derivatives are $\cos$ and $-\sin$; and every sinusoid has an acceleration
proportional and opposite to its displacement. In two or three dimensions
position, velocity and acceleration are vectors; the velocity is always
directed along the path, whereas the acceleration may point in any
direction, and in uniform circular motion it points towards the centre.
The Taylor series approximates a motion over a short interval with an
error that falls as a known power of the interval. The same definitions
apply to a system of atoms, whose motion is a path in $3N$ dimensions,
and velocities recovered from saved frames carry an error determined by
the interval between the frames. What determines the acceleration has so
far been left open, and Chapter~2 answers that question by introducing
force.
